\section{The Compound Exposure Investment Case}
\label{sec:investment-case}

\subsection{NPV Structure of Compound Investment}

The investment case for compound exposure corridors differs structurally from single-dimension investment cases. Define:

\begin{align}
\text{NPV}_{\text{compound}} &= \text{NPV}_{B1\text{-fix}} + \text{NPV}_{D1\text{-fix}} - \text{NPV}_{\text{overlap}} \label{eq:compound-npv}
\end{align}

where $\text{NPV}_{B1\text{-fix}}$ is the net present value of benefits from addressing structural isolation, $\text{NPV}_{D1\text{-fix}}$ is the net present value of benefits from addressing environmental vulnerability, and $\text{NPV}_{\text{overlap}}$ is the cost of the overlap --- the shared capital investment required by a project that addresses both simultaneously. In a single-dimension investment, the overlap term does not appear: the project addresses only one dimension. In a compound investment (such as a tunnel that both eliminates the closure-prone segment and creates a second path), the overlap term is negative in the NPV equation because a \textit{single} capital cost purchases \textit{both} benefit streams. This is the NPV superiority of compound investments.

\subsection{Donner Pass Freight Tunnel: The Benchmark Case}

The Donner Pass freight tunnel is the highest-NPV project in the compound exposure portfolio. The engineering concept follows the precedent of the original Central Pacific Railroad tunnel through the Sierra Nevada (at approximately the same elevation and geology) and more recent analysis of similar mountain tunnel investments \citep{FHWA_freight2023}.

\textbf{Cost estimate.} \$4.0 billion capital cost. This estimate is derived from analogous mountain tunnel construction: the California High Speed Rail Authority estimates \$4--5B per tunnel mile for similar Sierra Nevada geology (granite with fault zones); the proposed Donner freight tunnel is approximately 12--15 miles at similar geology. Railroad precedent (Union Pacific's ongoing tunnel feasibility work at Donner) suggests \$3.5--4.5B range. We use \$4.0B as the midpoint.

\textbf{B1 benefit.} Eliminating the single-point-of-failure at Donner Pass by providing a second interstate-grade crossing of the Sierra Nevada. The alternate path created by the tunnel reduces the B1 score from 8.3 to an estimated 4.0, creating genuine redundancy between I-80 and the tunnel. Annual benefit from avoided rerouting: 8,000 trucks/day $\times$ 365 days/year $\times$ closure probability (50 closures/year $\times$ mean 18 hours / 8,760 hours/year = 10.3\% of time) $\times$ avoided rerouting penalty (310 miles $\times$ \$4.50/mile + 5.5 hours $\times$ \$225/hr = \$2,625/truck affected). Annual rerouting benefit = 8,000 $\times$ 0.103 $\times$ 365 $\times$ 0.4 (fraction that would reroute rather than wait) $\times$ \$2,625 $\approx$ \$900M/year \citep{ATRI_costs2024}.

\textbf{D1 benefit.} The tunnel bypasses the weather-exposure segment entirely: I-80 over Donner Summit (elevation 7,227 feet) is replaced by a tunnel at sub-summit elevation, where snowfall and ice conditions are controlled. Annual benefit from avoided waiting costs: 8,000 trucks/day $\times$ 50 closures $\times$ 18 hours mean duration $\times$ \$225/hr $\times$ fraction stranded = 8,000 $\times$ 50 $\times$ 18 $\times$ \$225 $\times$ 0.6 (fraction stranded/waiting) $\approx$ \$700M/year \citep{ATRI_costs2024}.

\textbf{Combined NPV.} Total annual benefit \$1.6B (B1 + D1 components; no double-counting because rerouting and waiting are distinct behaviors). At a 7\% real discount rate over a 30-year infrastructure horizon:

\begin{equation}
\text{NPV}_{\text{tunnel}} = \sum_{t=1}^{30} \frac{\$1.6\text{B}}{(1.07)^t} - \$4.0\text{B} \approx \$19.8\text{B} - \$4.0\text{B} = \$15.8\text{B}
\end{equation}

Cost-benefit ratio: 5.75:1.

\textbf{Single-dimension comparison.} Table~\ref{tab:donner-comparison} shows the NPV comparison between the compound tunnel investment and the two single-dimension alternatives.

\begin{table}[h!]
\centering
\caption{Donner Pass: Compound vs.\ Single-Dimension Investment Comparison}
\label{tab:donner-comparison}
\begin{tabular}{lrrrr}
\toprule
\textbf{Investment} & \textbf{Cost (\$B)} & \textbf{Annual benefit (\$B)} & \textbf{30-yr NPV (\$B)} & \textbf{CBR} \\
\midrule
Freight tunnel (B1 + D1) & 4.0 & 1.60 & 15.8 & 5.75:1 \\
US-50 alternate route (B1 only) & 3.0 & 0.90 & 8.6 & 3.9:1 \\
Snowshed hardening (D1 only)  & 0.8 & 0.28 & 2.7 & 4.4:1 \\
\bottomrule
\end{tabular}
\end{table}

The tunnel dominates both single-dimension options by NPV. The US-50 alternate upgrades the existing highway to commercial-truck standard during winter closures, addressing the B1 isolation problem at lower capital cost than a tunnel --- but without reducing the 50 annual closures (which continue on I-80 itself). The snowshed hardening reduces D1 exposure by approximately 40\% (covering the four highest-closure segments), at low capital cost --- but without addressing the B1 isolation: when closures occur (even 40\% less often), there is still no adequate alternate. The tunnel simultaneously reduces D1 to near zero (no weather exposure in tunnel) and B1 from 8.3 to 4.0 (second crossing created), and this dual effect produces the superior NPV.

\subsection{Gulf Coast I-10: Hardening-Led Compound Investment}

Gulf Coast I-10 presents a different compound investment structure. The B1 score of 6.8 is moderate: I-20 provides a functional alternate through northern Louisiana, adding 180 miles but accessible to commercial freight. The D1 score of 8.4 is the highest for any T1 corridor in the corpus, driven by chronic flood exposure, subsidence, and hurricane surge probability.

The primary intervention is infrastructure hardening (D1 fix): roadway elevation at six lowest-lying segments, improved drainage structures, and I-10 causeway extension with hurricane-resistant pier design. Estimated cost: \$2.1B (Louisiana segment, Baton Rouge to Lake Charles). The B1 component is addressed through designated freight alternate status for I-20 (Shreveport corridor) and I-10B, at a cost of \$0.8B for interchange upgrades and load rating improvements. Total portfolio: \$2.9B.

Annual benefit: avoided closures from the 8 annual major inundation events (mean 36-hour closure, 12,000 trucks/day, \$225/hr operating penalty, \$1,800/truck affected) plus avoided I-20 detour cost for the fraction that reroutes \citep{ATRI_costs2024}. Combined annual benefit: approximately \$820M/year. 30-year NPV at 7\%: \$8.7B. Cost-benefit ratio: 4.0:1.

\subsection{Portfolio Comparison}

Table~\ref{tab:portfolio-npv} presents the investment case for the top four compound exposure corridors.

\begin{table}[h!]
\centering
\caption{Top-4 Compound Exposure Corridors: Investment Summary}
\label{tab:portfolio-npv}
\begin{tabular}{lrrrr}
\toprule
\textbf{Corridor} & \textbf{Cost (\$B)} & \textbf{Annual benefit (\$B)} & \textbf{30-yr NPV (\$B)} & \textbf{CBR} \\
\midrule
I-80 Donner Pass (tunnel)         & 4.0 & 1.60 & 15.8 & 5.75:1 \\
I-10 Gulf Coast LA (harden + alt) & 2.9 & 0.82 & 5.8  & 3.0:1  \\
I-90 Snoqualmie (snowshed + alt)  & 1.6 & 0.65 & 4.7  & 3.9:1  \\
I-5 Siskiyou Pass (harden)        & 1.1 & 0.48 & 3.2  & 3.9:1  \\
\midrule
\textbf{Portfolio total}           & \textbf{9.6} & \textbf{3.55} & \textbf{29.5} & \textbf{4.1:1} \\
\bottomrule
\end{tabular}
\end{table}

All four top-priority compound exposure investments exceed 3:1 cost-benefit ratio at a 7\% discount rate. The portfolio total of \$9.6B in capital investment produces \$29.5B in 30-year NPV --- a national-scale return that cannot be matched by equivalent investment in single-dimension resilience improvements across a larger set of lower-priority corridors.
